%% ma: L12-B13-0008
%% lop: 12
%% chuong: 4
%% bai: 13
%% dang: 12.13.1
%% loai: TLN
%% muc: VD
%% dap_an: "235"
%% nguon: "(NBV)-ĐỀ SỐ 2-MỖI NGÀY 1 ĐỀ THI 2025.pdf (D02, MỖI NGÀY 1 ĐỀ - Đề số 2), Phần III Câu 5"
%% kien_thuc: "Bài 13 (thể tích khối tròn xoay quanh trục Ox, phương trình elip)"
%% ghi_chu: "Trùng dạng 12.13.1 (thể tích khối tròn xoay quanh Ox), đã duyệt, mức VD. Đã thêm câu “quả dưa hấu có dạng khối tròn xoay ... quanh trục lớn” (đề gốc ngầm hiểu như vậy). Đề gốc có hình elip ghi mốc ±10 trên cả hai trục (không khớp trục lớn 28), bỏ hình. Đáp án gốc 235 đúng (kiểm bằng sympy: 234,57)"
%% trang_thai: chinh-thuc
\begin{ex}
Bổ ngang một quả dưa hấu, ta được thiết diện là hình elip có trục lớn $28$ cm, trục nhỏ $20$ cm. Quả dưa hấu có dạng khối tròn xoay tạo thành khi quay hình elip đó quanh trục lớn. Biết rằng ước tính cứ $500\ \mathrm{cm}^3$ dưa hấu sẽ làm được một cốc sinh tố giá $20\,000$ đồng. Từ quả dưa hấu trên có thể thu được $x$ (đơn vị: nghìn đồng) từ việc bán sinh tố. Tính giá trị của $x$ (làm tròn đến hàng đơn vị của nghìn đồng). Giả sử bề dày vỏ dưa không đáng kể.
\shortans{235}
\loigiai{Chọn hệ trục $Oxy$ có $Ox$ là trục lớn, elip có $a=14$, $b=10$: $\dfrac{x^2}{14^2}+\dfrac{y^2}{10^2}=1\Rightarrow y^2=100\left(1-\dfrac{x^2}{196}\right)$.
Thể tích quả dưa: $V=\pi\displaystyle\int_{-14}^{14}100\left(1-\dfrac{x^2}{196}\right)\mathrm dx=\dfrac{5600\pi}{3}\ (\mathrm{cm}^3)$.
Số tiền thu được: $\dfrac{5600\pi}{3}\cdot\dfrac1{500}\cdot20\,000\approx234\,572$ đồng $\approx235$ nghìn đồng.}
\end{ex}
